%% =====================================================================
%% formal.tex -- Sections 4 (Formal Semantics) and 5 (Verification Algorithm).
%% The algorithm of Figure fig:alg is transcribed from covenant/timing.py of the release.
%% =====================================================================
\providecommand{\Sc}[1]{\mathit{Sc}_{#1}}
\providecommand{\ASc}[2]{\mathit{ASc}_{#1}^{#2}}
\providecommand{\ESs}{\mathit{ES}^{\mathrm{src}}}
\providecommand{\ESd}{\mathit{ES}^{\mathrm{dst}}}
\providecommand{\AS}{\mathit{AS}}
\providecommand{\chainto}{\rightsquigarrow}
\providecommand{\mspecl}[1]{\ensuremath{\Rightarrow_{m}^{#1}}}
\providecommand{\ordl}[1]{\mathit{ord}_{#1}}
\providecommand{\Hpw}{\ensuremath{\Haz_C^{\mathit{pw}}}}
\providecommand{\Hep}{\ensuremath{\Haz_C^{\mathit{ep}}}}
\providecommand{\Hspec}{\ensuremath{\Haz_{\mathit{spec}}}}
\providecommand{\kw}[1]{\textbf{#1}}
\providecommand{\pc}[1]{\textsc{#1}}
\providecommand{\posof}{\pi}
\providecommand{\Hstrict}{\ensuremath{\Haz_{\mathit{spec}}^{\mathit{dom}}}}

%% ===================================================================
\section{Formal Semantics of the Recorded Frame}
\label{sec:model}

The guarantee of \sys is a chain of steps, from a run of the host program to a statement about that run
(Table~\ref{tab:chain}). Each step is of one of three kinds: a definition or theorem that we prove
(Appendix~\ref{app:proofs}), a named premise about capture or modelling (\S\ref{sec:premises}), or trusted code
that we test but do not prove (\S\ref{sec:tcb}). We keep the three apart.

\begin{table}[t]
\caption{From captured execution to stated guarantee. ``Proved'' means proved in Appendix~\ref{app:proofs} for the
algorithm of Figure~\ref{fig:alg}, relative to our transcription of the specification
(Definitions~\ref{def:events}--\ref{def:spec}). The transcription follows the Vulkan specification at Vulkan-Docs
tag \code{v1.4.365}, the revision of the headers our model driver and validation baseline use.}
\label{tab:chain}
\small
\begin{tabular}{@{}p{0.30\columnwidth}p{0.30\columnwidth}p{0.32\columnwidth}@{}}
\toprule
Step & Bridge & Kind \\
\midrule
run $\to$ trace $\tau$ & A1, A2-H; scope A3 & premise about capture \\
$\tau \to$ events $E(\tau)$ & A2, A-HOST, A-RT/EV/DS & premise about the model \\
$E(\tau) \to \Hspec(\tau)$ & Def.~\ref{def:events}--\ref{def:spec} & definition (our transcription) \\
algorithm $\to \Hspec$ & Thm.~\ref{thm:sound}, Lem.~\ref{lem:fast} & proved, under W1 (and W2 for Lem.~\ref{lem:fast}) \\
source code $\to$ algorithm & \S\ref{sec:tcb} & trusted, tested \\
no report $\to$ run is race-free & Cor.~\ref{cor:sched} & proved from the rows above \\
no premise warning $\to$ host domain & Prop.~\ref{prop:hostdom} & proved; removes the device-to-host part of A-HOST \\
recorded frames $\to$ later frames & Thm.~\ref{thm:window} & proved, conditional on P-PER, which the trace cannot establish \\
\bottomrule
\end{tabular}
\end{table}

\subsection{Events and synchronization}
\label{sec:events}

\begin{definition}[Trace, queues and events]\label{def:events}
A trace $\tau$ is the sequence of Vulkan calls of one run, in host order. From it \sys builds a set of queues
$Q = \{\mathit{host}\} \cup Q_{\mathit{dev}} \cup \{\mathit{present}{:}i\}_i$, with one presentation queue per
swapchain image, and for each queue $q$ a sequence of operations, appended in host order: the $k$-th operation of $q$
is at position $(q,k)$, and $\mathit{call}(o)$ is the host call that appended $o$. An operation is an access, a
barrier struct, a signal $\mathsf{Sig}(x,S)$ or a wait $\mathsf{Wait}(x,W)$ of a name $x$, or an event set/reset
marker. $S$ and $W$ are stage masks or $\top$ (all stages). A queue submission is $\mathsf{Sig}(\mathit{submit}_c)$ on
the host followed by $\mathsf{Wait}(\mathit{submit}_c,\top)$ at the head of the batch on its queue. A fence or idle
wait is a signal on a device queue and a $\top$ wait on the host. Swapchain acquire and present are encoded on the
presentation queues as in \S\ref{sec:impl}. A queue's operations after a wait whose name is never signalled are
dropped.

The events $E(\tau)$ are the accesses and the layout transitions; a transition $t_B$ sits at the position of its
barrier struct $B$. An event $a$ has a queue $q(a)$, a position $p(a)$, a resource $r(a)$ and a region
$R(a)$: a set of elements, which are the bytes of a buffer, or the aspect $\times$ layer slots of one mip level of an
image. It also has a kind $k(a) \in \{\mathsf{R},\mathsf{W}\}$ (transitions are $\mathsf{W}$), a stage $s(a)$ (a
pipeline stage, $\bot$ if unknown, or the pseudo-stage $\mathsf{LT}$ for a transition), an access flag $f(a)$ ($\bot$ if
unknown), and four exemption tags: its command, its render-pass scope, its subpass instance with a load/store
marker, and its atomic mode. A transition is \emph{explicit} if an image memory barrier performs it and
\emph{automatic} if a render pass does. An event may carry a \emph{position token} $\posof(a)$
(\S\ref{sec:fblocal}); $\posof(a) = \bot$ otherwise. The frames of $\tau$ are the segments of the host's call sequence that end at each
\code{vkQueuePresentKHR}; $\mathit{fr}(o)$ is the frame that contains $\mathit{call}(o)$.
\end{definition}

\paragraph{Scopes}
For a stage mask $M$, $\ESs(M)$ and $\ESd(M)$ are its source and destination \emph{execution} stage sets. They add to
each named stage the logically earlier (source) or later (destination) stages of every synchronization pipeline
that contains it, as listed in the Vulkan registry. \code{ALL\_COMMANDS}, \code{BOTTOM\_OF\_PIPE} (source) and
\code{TOP\_OF\_PIPE} (destination) give $\top$; \code{TOP\_OF\_PIPE} (source) and \code{BOTTOM\_OF\_PIPE}
(destination) give $\emptyset$. \code{ALL\_GRAPHICS} gives the graphics and mesh stages, \code{TRANSFER} also gives
\code{COPY}, \code{BLIT}, \code{CLEAR} and \code{RESOLVE}, and the acceleration-structure build stage also gives the
copy stage. $\AS(M)$, the \emph{access} stage set, applies only these group expansions; it adds no earlier or later
stages. $\top$ contains every stage except $\mathsf{LT}$, which is in a set only if it is named, and $\bot$ is in a set
only if the set is $\top$. The two fragment-test stages are interchangeable in access sets (decision M1, as in
validation). $f \in_k A$ is access-mask membership for an access of kind $k$; \code{MEMORY\_READ} and
\code{MEMORY\_WRITE} contain every read or write, including one whose flag is $\bot$, which no other mask contains,
and \code{SHADER\_READ} does not contain \code{UNIFORM\_READ}.

\begin{definition}[Dependencies, chains, ordering and races]\label{def:spec}
A \emph{dependency} $D$ is a barrier struct $B=(S_1,A_1,S_2,A_2,r_B,R_B)$, where $r_B = \ast$ for a global
barrier; a signal; or a wait. Each $D$ has a position $p(D)$ on a queue $q(D)$. A barrier struct also has a call
identity, and an event wait has a cut $c_B$: the position of the latest set of its event before it, or $-\infty$ if
there is none or a reset followed it. Each dependency has a first and a second stage set: $\ESs(S_1)$ and $\ESd(S_2)$
for a barrier, $\ESs(S)$ for a signal and $\ESd(W)$ for a wait (empty on the side that does not apply). Host signals,
host waits and submission waits have $\top$.

\emph{Locality.} A barrier struct or subpass dependency $D$ also has a locality $\lambda(D)$: $\mathsf{loc}$ when its
dependency flags hold \code{BY\_REGION}, or \code{VIEW\_LOCAL} with view offset~0; $\mathsf{nopos}$ for
\code{VIEW\_LOCAL} with a non-zero view offset; and $\mathsf{glob}$ otherwise. Signals and waits are $\mathsf{glob}$.
$F$ is the set of framebuffer-space stages (fragment shader, early and late fragment tests, color attachment output).
An event $a$ is \emph{located} if $s(a) \in F \cup \{\bot, \mathsf{LT}\}$, that is, if it may be a framebuffer-space
operation.

\emph{Synchronization scopes.} An event $a$ with $q(a) = q(D)$ is in $\Sc{1}(D)$ if $p(a) < p(D)$, $s(a)$ is in
$D$'s first stage set, $a$ is not in another struct of $D$'s call, and $p(a) < c_D$ when $D$ is an event wait. It
is in $\Sc{2}(D)$ if $p(a) > p(D)$ and $s(a)$ is in $D$'s second stage set. Transitions are the exception: $t_B$ is
in $\Sc{1}(B)$ and $\Sc{2}(B)$, in $\Sc{1}$ of a later signal with $S = \top$, and in $\Sc{2}$ of an earlier wait
with $W = \top$ or an earlier submission wait, and in no other scope. This is membership in scopes, not in stage
sets: where the ordering below tests the stage of $t_B$ against a dependency whose scope this rule puts $t_B$ in,
the test holds, and for the link classes $t_B$ counts as framebuffer-space ($\mathsf{G}$ and $\mathsf{L}$ links
hold it, $\mathsf{N}$ links do not). No API stage mask contains $\mathsf{LT}$, so two dependencies never chain
through a transition (rule R-LT). So a write in $B$'s first access scope is ordered before $t_B$ through the link
$(B,\mathsf{G})$, and $t_B$ before an access in $B$'s second access scope.

\emph{Access scopes,} for an element $\ell$: $a \in \ASc{1}{\ell}(D)$ if $a \in \Sc{1}(D)$ and, when $D$ is a barrier
struct, $s(a) \in \AS(S_1)$, $f(a) \in_{k(a)} A_1$, and $r_B = \ast$ or ($r_B = r(a)$ and $\ell \in R_B$).
$\ASc{2}{\ell}(D)$ is defined in the same way from $S_2$ and $A_2$. A signal's access scopes are its whole
synchronization scopes, and so are a wait's. A transition is in both access scopes of its own barrier.

\emph{Links.} A \emph{link} $(D,c)$ is a dependency with a class $c \in \{\mathsf{G},\mathsf{N},\mathsf{L}\}$ that
restricts its second stage set: to all of it ($\mathsf{G}$), to its stages outside $F$ ($\mathsf{N}$), or to its
stages in $F$ ($\mathsf{L}$). A link is \emph{entered locally} when $\lambda(D) \ne \mathsf{glob}$ and it is entered
from a located event or from an $\mathsf{L}$ link. A link entered locally has class $\mathsf{N}$ or $\mathsf{L}$
($\mathsf{L}$ only if $\lambda(D) = \mathsf{loc}$); every other link has class $\mathsf{G}$.

\emph{Chaining.} $(D,c) \chainto (D',c')$ if $D$ and $D'$ are on one queue, $p(D) < p(D')$, $D$'s second stage set
restricted by $c$ meets $D'$'s first stage set in a stage other than $\mathsf{LT}$ (or in $\mathsf{LT}$, when
$D$ and $D'$ are one struct of one call: the two halves of a split transition, \S\ref{app:elide}), they are not two
structs of one call, $p(D) < c_{D'}$ when $D'$ is an event wait, and $c'$ follows the rule above; or if $D$ is $\mathsf{Sig}(x,\cdot)$, $D'$ is $\mathsf{Wait}(x,\cdot)$
and $c = c' = \mathsf{G}$. On the host, every earlier dependency chains into every later one.

\emph{Ordering.} $a \xspec b$ if $q(a) = q(b) = \mathit{host}$ and $p(a) < p(b)$, or if there is a chain $(D_1,c_1)
\chainto \dots \chainto (D_n,c_n)$ ($n \ge 1$) with $a \in \Sc{1}(D_1)$, $c_1$ given by the rule above with $D_1$
entered from $a$, $b \in \Sc{2}(D_n)$ with $s(b)$ in $D_n$'s second stage set restricted by $c_n$, and
$\posof(a) = \posof(b) \ne \bot$ if $c_n = \mathsf{L}$. $a \mspecl{\ell} b$ holds if the same is true with $a \in
\ASc{1}{\ell}(D_1)$ and $b \in \ASc{2}{\ell}(D_n)$; host program order is also a memory order, and a host wait is
treated as one (premise A-HOST, \S\ref{sec:hostdom}). The intermediate dependencies need only chain.

\emph{Races.} $a$ and $b$ \emph{conflict} if $r(a) = r(b)$, $R(a) \cap R(b) \ne \emptyset$, one of them writes, and
they are not exempt. Five kinds of pair are exempt: (E1)~two accesses of one command; (E2)~attachment accesses of one
render-pass scope on one queue (rasterization order); (E3)~a load or store operation and an access of its subpass
instance; (E4)~two accesses on one queue with compatible atomic modes, that is, wherever the bytes they may touch
overlap, both are atomic operations whose memory scope holds every invocation of the device (\code{Device},
\code{QueueFamily} or \code{CrossDevice}), or both read (bytes are buffer locations: a block member's offset plus
the descriptor's offset; an access whose member is not known may touch the whole range); (E5)~two transitions on one queue, at least one of them explicit (rule R-LT).
A store through an index taken from an atomic is an ordinary write: the atomic need not hand out distinct indices
(\code{atomicAdd(c, 0)} gives every caller the same one).
They \emph{race at} $\ell \in R(a) \cap R(b)$ if neither $\ordl{\ell}(a,b)$ nor $\ordl{\ell}(b,a)$ holds, where
$\ordl{\ell}(x,y)$ is $x \mspecl{\ell} y$ if $k(x) = \mathsf{W}$ and $x \xspec y$ otherwise. We write $a \spec b$ when
$\ordl{\ell}(a,b)$ holds for every $\ell \in R(a) \cap R(b)$. $\Hspec(\tau)$ is the set of pairs that race at some
element.
\end{definition}

\paragraph{Rule R-LT}
Layout transitions ``performed via image memory barriers execute in their entirety in submission order, relative to
other image layout transitions submitted to the same queue, including those performed by render passes'', so a pair of
transitions on one queue of which at least one is by an image memory barrier is exempt (E5). Two automatic
transitions of render passes are outside this sentence and are checked like any other pair (the matrix of transition
pairs, with the specification text and a regression test for each cell, is in \supp{B}). The order is not
transitive: ``an application must not rely on ordering of image layout transitions to influence ordering of other
commands''. R-LT therefore adds no chain, and a write ordered before the first transition but not before the second
still races the second. This wording entered the specification in v1.4.326; earlier revisions gave transitions an
implicit dependency, so R-LT is the newer, weaker rule. The same paragraph adds that the ordering ``only applies if the
implementation performs actual read/write operations during the transition''; R-LT is consistent with it, since a
transition that performs no read or write leaves nothing for the exempt pair to race on, and R-LT never orders other
commands through a transition.

\paragraph{Framebuffer-local dependencies}
\label{sec:fblocal}
A dependency with \code{BY\_REGION} is \emph{framebuffer-local} for a pair of framebuffer-space operations: it orders
them only where their framebuffer regions overlap, and \code{VIEW\_LOCAL} adds the same restriction per view. A pair
in which either operation is in a stage outside framebuffer space is ordered as by any other dependency, and a chain
with one ordinary link is ordinary. Definition~\ref{def:spec} encodes this with the link classes: $\mathsf{N}$ carries
the ordinary part of a local dependency entered from a located event, $\mathsf{L}$ its framebuffer-space part, and any
ordinary link returns the chain to $\mathsf{G}$. An $\mathsf{L}$ ending orders only an access at the source's own
position, which the trace translation records as a token: for an attachment read or write, a load, store or resolve
operation, an automatic transition of an attachment, or a single-sample input-attachment read, $\posof(a)$ is the image
with the base mip level and base layer of the view through which it is accessed. Every other access has
$\posof(a) = \bot$: a storage, sampled or uniform access of one fragment can reach any byte, so no local dependency
orders it.

\begin{lemma}[Position tokens]\label{lem:tokens}
Let $\posof(a) = \posof(b) \ne \bot$. Then every texel that $a$ and $b$ both access is accessed by each of them at the
same framebuffer coordinate and view.
\end{lemma}
\begin{proof}
A framebuffer-space operation with a token accesses, at framebuffer coordinate $(x,y)$ of view $v$, the texel
$(x,y)$ of layer $\mathit{base}+v$ of mip $\mathit{base}$ of its image, and no other texel; a single-sample input
attachment read is restricted to its own pixel by the specification. Two operations whose views have the same image,
base mip and base layer therefore map each common texel to the same coordinate and view.
\end{proof}

So an $\mathsf{L}$ chain from $a$ orders $b$ at every common texel, and treating \code{nopos} links as $\mathsf{N}$
only, and tokenless accesses as unordered by $\mathsf{L}$ links, can only add races. Traces recorded without subpass
dependency flags are read with every explicit subpass dependency local, the weaker reading. Theorem~\ref{thm:sound}
and Lemma~\ref{lem:fast} are proved for this definition (Appendix~\ref{app:sound}), and the executable reference
implements it.

\paragraph{The verified configuration}
Definition~\ref{def:spec} follows the specification on layout transitions from \code{UNDEFINED}: a write before such a
transition must be made available to it ($\ordl{\ell}$ uses $\mspecl{\ell}$ when the earlier access writes), although the
transition discards the contents. The specification requires that all memory writes be made available before a
layout transition executes, and \code{UNDEFINED} only permits the old contents to be discarded. The \emph{verified
configuration} is \sys's default: spec-strict transitions, conservative host capture (\S\ref{sec:capture}) and the
weaker reading of unrecorded subpass flags. Two opt-in switches leave it, a heuristic that accepts an execution
dependency before a transition from \code{UNDEFINED} (relaxation R1) and byte-difference host capture; each report
records its configuration, and every theorem in this paper is stated for the verified one.

%% ===================================================================
\subsection{Premises}
\label{sec:premises}

Table~\ref{tab:assume} lists the premises. They are of three kinds, plus the periodicity premises of
\S\ref{sec:window}. \emph{Capture} premises (A1, A2-H) relate the run
to the trace, and A3 states the scope. \emph{Model} premises (A2, A-HOST and the rows for ray tracing, events and
dynamic state) relate the trace to the events. \emph{Well-formedness} premises (W1, W2) hold by construction of the
trace translation. The checker's own correctness is not a premise. Section~\ref{sec:tcb} states which code is trusted and how it is
tested.

\begin{table}[tp]
\caption{Premises of the theorems, and their status. TLAS/BLAS: top-/bottom-level acceleration structure; SBT:
shader binding table; BDA: buffer device address.}
\label{tab:assume}
\footnotesize
\begin{tabular}{@{}p{0.27\columnwidth}p{0.68\columnwidth}@{}}
\toprule
Premise & Status \\
\midrule
\multicolumn{2}{@{}l}{\emph{Capture}} \\
A1 trace faithfulness & by construction of the recording driver: every call of the run with its arguments, in one
host order. A command the model does not interpret is reported (\code{timing.unmodelled-command}); no claim is made for it \\
A2-H conservative host capture & page protection: a page written between two calls is charged to every binding it
overlaps, at that call boundary (\S\ref{sec:capture}). The opt-in byte-difference capture is outside A2-H \\
A3 one host path (scope) & only the recorded path and one interleaving of host threads \\
\midrule
\multicolumn{2}{@{}l}{\emph{Model}} \\
A2 device accesses over-approximated & an event per access of a recorded command, with its stage and flag (or $\bot$)
and a region that contains it; exemption tags never claimed falsely. Byte intervals; aspect $\times$ layer per mip;
static binding use; every written element of a bound descriptor array, except under a holding index guard
(Lemma~\ref{lem:guard}). Aliasing of distinct resources is not modelled \\
A-HOST host-domain operations & host-to-device: a submission makes earlier host writes available and visible to the
device. This holds for host-coherent memory, the only host-visible memory the model driver offers, so it holds for
the recorded run; for non-coherent memory on another device it needs the application's
\code{vkFlushMappedMemoryRanges}. Device-to-host: a host wait (fence, timeline value, query with \code{WAIT}, idle)
orders the waited scope before later host events (specified), and Definition~\ref{def:spec} also treats it as a memory
order, which the specification does not grant. \sys reports every pair that relies on the latter as a premise warning;
with none reported, A-HOST is used only for host-to-device (Proposition~\ref{prop:hostdom}). Host reads are not events,
so no claim on host readback \\
A-RT1 TLAS references & read at submission from host-visible instance buffers; if unresolvable, every live BLAS \\
A-RT2 device addresses & resolved at record time against live buffers; an unmatched address is reported \\
A-RT3 shader BDA accesses & not tracked; the pipeline is listed as ``untracked'' \\
A-RT4 query results & not memory (AS properties, compacted sizes, timestamps) \\
A-RT5 \code{vkCmdTraceRays}\-\code{Indirect2KHR} & SBT in GPU memory, unknown; reported \\
A-RT6 host AS builds/copies & a host write of the destination; inputs not tracked \\
A-EV1 events & a wait is a barrier whose first scope is cut at the latest set of the event on its queue; a host
\code{vkSetEvent} orders nothing on the device (reported) \\
A-DS1 dynamic state & whitelisted state folded into the over-approximation (all bound vertex buffers read, all color
attachments written) \\
\midrule
\multicolumn{2}{@{}l}{\emph{Well-formedness} (by construction of the trace translation)} \\
W1 unique signals & each name is signalled at most once (numbered binary instances, timeline values, fence
generations) \\
W2 contiguity & the accesses of one command lie on one queue with nothing between them; no wait, signal or access of
another scope (instance) lies between two accesses of one render-pass scope (subpass instance); the structs of one
barrier call are consecutive (Lemma~\ref{lem:fast} only) \\
\midrule
P-MONO, P-WSI, P-PER & Theorem~\ref{thm:window} only (\S\ref{sec:window}); the first two and periodicity inside
$\tau$ can be checked on $\tau$, continuation beyond it cannot \\
\bottomrule
\end{tabular}
\end{table}

%% ===================================================================
\subsection{Soundness}

Let $\Hpw(\tau)$ be the set of pairs that the pairwise mode of Figure~\ref{fig:alg} reports, and $\Hep(\tau)$ the set
that the epoch mode reports.

\begin{theorem}[Soundness of the checking algorithm]\label{thm:sound}
Let $\tau$ satisfy W1, and run the algorithm of Figure~\ref{fig:alg} in the verified configuration.
If \pc{Ordered}$(a,b)$ returns \emph{ok} for conflicting events $a, b$, then $a \spec b$. Hence
$\Hspec(\tau) \subseteq \Hpw(\tau)$.
\end{theorem}

The proof (Appendix~\ref{app:sound}) follows the algorithm's own structure, in five lemmas. (i)~The order in which the
algorithm processes operations extends program order and the pairing of signals with waits. (ii)~The barrier chain
state of an event reaches exactly the links, with their classes, that a chain of Definition~\ref{def:spec} reaches
from it; in memory mode it reaches only links whose first dependency covers the whole event region. (iii)~A same-queue answer of
\emph{ok} exhibits such a chain. (iv)~An accepted relay entry of a wait record (an entry at or after $p(e)$ that passes its stage filter)
means that a chain leads from $e$ into that wait; this lemma needs W1. (v)~Each way in which a wait is accepted puts
the later event in a scope of a chain that starts at that wait. Finally, the pairwise mode presents every
earlier conflicting event as a candidate. A race is therefore never answered \emph{ok}, and it is reported unless it
is exempt, and exempt pairs are not races.

\begin{corollary}[What an accepted trace guarantees]\label{cor:sched}
Assume also A1, A2, A2-H, A-HOST and the modelling rows of Table~\ref{tab:assume}. If $\Hspec(\tau) = \emptyset$
(in particular, if the pairwise mode reports nothing, or if the epoch mode reports nothing under the premises of
Lemma~\ref{lem:fast}), then on the recorded host path every two conflicting memory accesses that the model covers are
ordered as Definition~\ref{def:spec} requires. Hence no device schedule consistent with $\tau$'s synchronization lets
them race. If, moreover, no premise warning is reported, this holds with host waits given only their specified
execution dependency (Proposition~\ref{prop:hostdom}).
\end{corollary}

The bridge from the run to $E(\tau)$ is an abstraction lemma (Appendix~\ref{app:abstraction}). A2 and A2-H give each
covered access $\alpha$ of the run a representative event $a$ whose scopes are subsets of $\alpha$'s. An ordering
of the representatives therefore implies the same ordering of the accesses, and every race of covered accesses is a
race of their representatives.

\subsection{Host-domain operations}
\label{sec:hostdom}
A device write is visible to a later host write of the same memory only after a \emph{domain operation}: a memory
dependency whose destination access mask holds \code{HOST\_READ} or \code{HOST\_WRITE}. A fence or a host wait alone
gives an execution dependency; the specification states that it does not make device results visible to the host.
Definition~\ref{def:spec} treats a host wait as a memory order, so on its own it would accept a device write followed,
after a fence wait, by a host write with no such barrier. We therefore separate that case.

\begin{definition}[Host-domain ordering]\label{def:hostdom}
A device write $e$ is \emph{host-domain ordered} before a host write $u$ if there is a barrier struct $B$ with $e
\mspecl{\ell} B$ for every $\ell \in R(e) \cap R(u)$ (a chain from $e$ whose last link is $B$), $\mathsf{HOST}$ in $B$'s
second stage and access stage sets, \code{HOST\_READ}, \code{HOST\_WRITE} or \code{MEMORY\_*} in its second access
mask, and a chain from $B$ whose last link holds $u$ in its second scope. $\Hstrict(\tau)$ is $\Hspec(\tau)$ plus every pair of a device write and a later host write
that conflict and are not host-domain ordered.
\end{definition}

\sys checks this for every pair it accepts in which a device write precedes a host write: it looks for such a $B$ on
the write's queue, reached by the write's memory chain and recorded before the latest position of that queue that a
host wait holds through signals of every stage. When there is none, it reports the pair as a premise warning
(\code{timing.premise-host-domain}). The test is conservative: a domain operation on another queue is not credited.

\begin{proposition}[Host domain]\label{prop:hostdom}
Under the premises of Theorem~\ref{thm:sound}, if the pairwise mode, or the epoch mode under W2, reports no error and
no premise warning, then $\Hstrict(\tau) = \emptyset$.
\end{proposition}

The proof (Appendix~\ref{app:hostdom}) shows that an accepted pair without a warning has the barrier the definition
asks for, and that in epoch mode a write that stands in for a dropped device write carries the domain operation over to
it, as in Lemma~\ref{lem:fast}. In the opposite direction no premise is left for the recorded run: the model driver
offers only host-coherent host-visible memory, where a submission performs the host-to-device domain operation. On a
device whose application chooses non-coherent memory, the host path differs (A3) and its flushes become a premise.

\subsection{Descriptor arrays}
\label{sec:guards}
A2 charges a draw with a read of every written element of each descriptor array its shaders use statically, which is
sound but imprecise when a shader selects one element per instance. \sys narrows the read only when it can show which
elements a draw can select. An \emph{index guard} $g$ of a direct, non-indexed draw $d$ holds when: (C1) every index
into the array is a constant or a chain of at most four constant shifts and masks of one flat 32-bit integer vertex
attribute; (C2) the pipeline has only vertex and fragment stages, no dynamic vertex-input state and no instance
divisor; (C3--C4) every attribute word $d$ fetches lies in host-written pages with recorded contents, rebuilt as of
$d$'s submission; (C5) no device command writes those bytes, and no image is bound on them; (C6) no host write to them
races $d$'s vertex-input read, and none recorded after the submission is ordered before it; (C7) every decoded index is
in range; (C8) the array holds sampled images. The decoded set $I(d)$ is the image of the decoded words under the
index chain. Under a holding guard, elements outside $I(d)$ are charged only if the guard fails.

\begin{lemma}[Index guards]\label{lem:guard}
Let $g$ hold for draw $d$ in the final check. Then in every device schedule consistent with $\tau$'s synchronization,
every invocation of $d$ indexes the array only with elements of $I(d)$, so the narrowed events still represent every
access of the run (A2).
\end{lemma}

The proof (Appendix~\ref{app:guard}) has three steps: the vertex fetch reads the rebuilt bytes in every schedule,
because no device command writes them (C5), every host write to them is ordered before the draw at or before its
submission or makes the guard fail (C6), and a race with any of them would be reported by Theorem~\ref{thm:sound}; the
shader maps those bytes into $I(d)$ (C1, C2, C7); and the events dropped therefore stand for no access. The guard is
not self-supporting: C5 and C6 mention only events on the attribute bytes, which are buffer bytes, and the
happens-before order, while the events a guard removes are image reads (C8), and accesses create no dependencies. So
the set of failing guards is the same whether or not the conditional reads are present; the checker computes it by
re-running with failed guards widened until no new guard fails. Any other index source (push constants, uniforms,
storage buffers, indirect or indexed draws) keeps the full charge.

\subsection{Guarantee eligibility}
\label{sec:eligible}
Every check reports three facts separately. \emph{Checking complete}: every queue ran to its end (no wait was left
without a signal). \emph{Reporting complete}: report grouping finished within its budget; when it does not, every
diagnostic is still reported, as a site group without a repair, and the result is marked incomplete (grouping never
removes a diagnostic, so this does not affect the guarantee). \emph{Guarantee-eligible}: the verified configuration,
checking complete, no detected capture or coverage violation (an unmodelled command, an untracked mapping, a capture
warning, host writes from another thread, a pipeline whose shader code is missing or cannot be read, or that reaches
descriptor memory through a value the shader analysis cannot attribute to a descriptor), and no premise
warning. The configuration is resolved once per check, from the environment and any explicit argument, and the
same values drive the check and the reported eligibility, so a check run with an opt-in heuristic is never reported
as verified. For an eligible trace, an empty error report is an
instance of Corollary~\ref{cor:sched} and Proposition~\ref{prop:hostdom}; P2 and the model premises remain assumptions.
\S\ref{sec:eval} reports eligibility per evaluated trace.

\paragraph{What the theorem covers, and what it does not}
\label{sec:cover}
\emph{Proved:} Theorem~\ref{thm:sound} and Lemma~\ref{lem:fast} relate the algorithm to
Definitions~\ref{def:events}--\ref{def:spec}. They are proved by hand, not mechanized, and relative to our
transcription of the specification, whose fidelity the differential tests of \S\ref{sec:eval-sem} test against
synchronization validation. \emph{Assumed about capture:}
\begin{itemize}
\item \emph{trace faithfulness} (A1): the trace holds the calls of this run, in host order;
\item \emph{conservative host capture} (A2-H): for every CPU store $\sigma$ to mapped memory, the trace has a host
  write event $w$ whose region contains $\sigma$'s bytes on every resource bound there, and every event ordered
  before $w$ (after $w$) is ordered before (after) $\sigma$. Page protection gives this: the first store to a page
  after a call faults and marks the page, the driver re-protects every page at the next call, and the write is
  placed between those two calls, where the host queue has no other operation. It does not depend on whether the
  store changed a value. The byte-difference mode charges only changed bytes, and is therefore a heuristic outside
  the corollary; so is the page-straddle rule that charges a shared page only to recently bound buffers;
\item \emph{the recorded host path only} (A3): no other path or host-thread interleaving.
\end{itemize}
\emph{Assumed about the model:} over-approximated device accesses (A2), the host-domain operations (A-HOST) and the
modelling rows of Table~\ref{tab:assume}. \emph{Not covered:} host paths the run did not take, including other
driver responses such as image counts, \code{VK\_SUBOPTIMAL\_KHR} results and swapchain recreation; other
interleavings of a multithreaded host; host \emph{reads} of mapped memory; accesses through aliased memory of distinct
resources; shader device-address accesses (A-RT3), query results (A-RT4), the inputs of host acceleration-structure
builds (A-RT6) and unmodelled commands; memory ordering inside a shader; and semantic conventions. The protocol checks
of \S\ref{sec:protocol} are separate guarantees, not consequences of Theorem~\ref{thm:sound}.

%% ===================================================================
\subsection{Precision}
\label{sec:precision}

An exact access abstraction does not, by itself, make the checker exact. Exact accesses fix $E(\tau)$, but they do
not make the decision procedure exact. The algorithm decides whether a
dependency covers an access with one region test (``the barrier's region contains the whole access region''), while
Definition~\ref{def:spec} decides element by element. For example, let a write $a$ cover bytes $[0,256)$, let one
barrier call carry two buffer structs for $[0,128)$ and $[128,256)$, and let a read $b$ of $[0,256)$ follow. Every
byte is ordered, so the pair is not a race. Neither struct contains all of $R(a)$, however, so the pairwise mode
reports it. What does hold is the following.

\begin{proposition}[Conditional precision of the pairwise mode]\label{thm:complete}
Let $\tau$ satisfy the premises of Theorem~\ref{thm:sound}. Assume further that (i)~no signal has a stage mask other
than $\top$, and (ii)~$\tau$ is \emph{region-uniform}: for every barrier struct $B$ with a region on resource $r$ and
every event $a$ on $r$, either $R_B \supseteq R(a)$ or $R_B \cap R(a) = \emptyset$. Then $\Hpw(\tau) = \Hspec(\tau)$.
\end{proposition}

Condition (ii) is the missing condition, and it can be met by construction. Splitting every access at all barrier
region bounds on its resource yields a region-uniform trace $\tau^\sharp$ with the same races (the parts of one access
are disjoint, so they never conflict, and Definition~\ref{def:spec} is decided per element; Appendix~\ref{app:precision}).
The pairwise mode on $\tau^\sharp$ is exact under (i). The checker does not split, so it may report
pairs that are not races when barrier regions cut accesses (the generator setting with partial region coverage,
\S\ref{sec:eval-sem}).
Without (i), the relay of a masked signal keeps a bounded number of stage-filtered entries per queue
(Figure~\ref{fig:alg}), and dropping one only removes orderings, so the checker can be stricter there.
The proposition is exact relative to $E(\tau)$. Relative to the \emph{run}, a report can also come from the
abstraction: an over-approximated region, an unknown stage or flag (which lies only in scopes of every stage or every
access), a captured host page, or a choice of the transcription that is stricter than the specification
(\S\ref{sec:eval-wholeframe}). Exact events remove the first three, but not a region cut, which is why an exact access
abstraction alone does not make the checker exact. The epoch mode reports a subset of $\Hpw$
(Lemma~\ref{lem:modes}), so it can miss pairs, but it flags at least one event of every race (Lemma~\ref{lem:fast}).

%% ===================================================================
\subsection{Periodic executions}
\label{sec:window}

A trace covers finitely many frames. A program usually runs longer, so we ask when a finite window decides the race
freedom of every later frame. The answer is conditional, because the future is not in the trace.

\begin{definition}[Periodic continuation]\label{def:periodic}
$\tau$ is \emph{periodic after warm-up $w$ with period $p$} if there is a shift $\sigma$ from the operations of frames
$\ge w$ onto those of frames $\ge w+p$. The shift preserves queues (up to a permutation of the swapchain images),
program order, all attributes up to a renaming of resources and synchronization objects, and the pairing of waits with
signals, and it maps frame $i$ onto frame $i+p$. Its \emph{periodic continuation} $\tau^\infty$ is the unique infinite
event structure whose first frames are those of $\tau$ and on which $\sigma$ extends to all frames $\ge w$.
\end{definition}

The premises are these. \textbf{P-MONO}: every wait is paired with a signal appended earlier in host order (no
wait-before-signal). \textbf{P-WSI}: every device access to a swapchain image is recorded in a command submitted
after an acquire of that image and before its next present, and no other event touches a swapchain image.
\textbf{P-PER}: the run continues as $\tau^\infty$. The parameters are properties of $\tau^\infty$; because $\sigma$
preserves them, they can be measured on the first $w+p+D$ frames of $\tau$. $N$ is the least number such
that, before the first submission or host access of each frame $j \ge w$, the host has waited on a $\top$ signal of
every device queue that follows all that queue's operations of frames $\le j - N$. $G$ is the largest frame distance
between a device access $b$ to a swapchain image and the acquire it is held under or the present that ends it.

\begin{theorem}[Window, conditional on periodic continuation]\label{thm:window}
Let $\tau$ be periodic after $w$ with period $p$, satisfy P-MONO, P-WSI and the premises of Theorem~\ref{thm:sound},
and let $D = \max(N, G+1)$ and $M = w + p + D - 1$. Every race of $\tau^\infty$ has a translate of a witness race in
the first $M$ frames. Hence, if $\tau$ has at least $M$ frames, $\Hspec(\tau^\infty) = \emptyset$ if and only if the
first $M$ frames of $\tau$ have no race. Under P-PER, a checker-clean window therefore implies race freedom of the
whole run (Corollary~\ref{cor:sched}).
\end{theorem}

The derivation (Appendix~\ref{app:window}) uses three facts. First, by P-MONO every chain runs forward in host order,
so the ordering of two events depends only on the operations appended between them, and a prefix of $\tau^\infty$ has
the same races as $\tau^\infty$ on its events. Second, two events that are not presentation-engine events and lie $N$
or more frames apart are ordered through the host's fence or idle wait. Third, presentation-engine accesses are
\emph{not} bounded by any frame fence. A fence proves that a submission has completed, never that the present waiting
on its semaphore, or the engine's read of the image, has. They are bounded through the acquire of the same image
instead. Under P-WSI, if a device access $b$ races any presentation-engine read of its image, then it also races
either the read that ends at the acquire that $b$ is held under, or the read of the present that ends it, and both lie
within $G$ frames of $b$. $G$ is 0 when every frame acquires, renders to and presents its own image.

\paragraph{What the trace can and cannot establish}
P-MONO, P-WSI, $N$, $G$ and periodicity within $\tau$ are properties of the trace, and can be checked on it.
P-PER is a claim about host behaviour after the trace ends, which input, resizing or streaming can change, and no
finite trace establishes it. Theorem~\ref{thm:window} therefore does not by itself license applying a barrier plan to
frames that were not traced. Section~\ref{sec:elide} restricts the transformation guarantee to the recorded execution,
and makes no claim about reusing a plan beyond it. The window analysis of the checker
computes $N$ from device queues only, compares frame signatures that abstract which signal each wait consumes, and
checks $w + N + 2$ frames. That bound reaches $M$ only when $p + \max(N, G+1) \le N + 3$, for example when $p \le 3$
and $G < N$, and its $N$ ignores the presentation queues. We therefore report it as a measurement on the traces and
not as an instance of the theorem (\S\ref{sec:eval-sem}).

%% ===================================================================
\subsection{Trusted computing base}
\label{sec:tcb}

Theorem~\ref{thm:sound} and Lemma~\ref{lem:fast} are proved for the algorithm of Figure~\ref{fig:alg}, not for its
Python source. We do not assume that ``the checker implements the specification relation''. Instead we state an
explicit trusted computing base: the code whose agreement with the model we trust, and how each
part is tested (Table~\ref{tab:tcb}). Everything else (report grouping and repair synthesis, diagnostic text, the
structure checks, the window analysis, the optimizer's search) is outside it. Their errors can cost precision or
readability, but they cannot turn a race into an accepted trace, because every verdict passes through T4.

\begin{table}[t]
\caption{Trusted computing base and how each part is tested.}
\label{tab:tcb}
\small
\begin{tabular}{@{}p{0.30\columnwidth}p{0.64\columnwidth}@{}}
\toprule
Trusted code & Tested by \\
\midrule
T1 recording driver (\code{covvk.cpp}): calls, arguments, host-write pages (A1, A2-H) & the same binaries under
synchronization validation in the whole-frame (1,251 programs) and ray-tracing/mesh (13) differential tests; capture
tests \\
T2 trace translation (\code{host\_trace}, \code{ranges}, SPIR-V access analysis): events, regions, exemption tags,
W1, W2 & unit tests, including assembled SPIR-V that varies helper calls, descriptor offsets and atomic against
plain accesses; the two whole-frame differential tests; 114 \bench reproductions. The executable reference consumes
T2's events, so it does not test T2 \\
T3 stage and access algebra (pipeline tables, scope expansion) & regression unit tests; an executable reference
with independently written tables; 750 barrier mutants under validation \\
T4 checker core (run loop, \pc{Ordered}, chain state, \pc{Access}) & executable reference agreement
(\S\ref{sec:eval-sem}); property tests of $\Hspec \subseteq \Hpw$, $\Hep \subseteq \Hpw$ and criterion E on random
frames; the differential tests above \\
\bottomrule
\end{tabular}
\end{table}

%% ===================================================================
\section{Verification Algorithm}
\label{sec:algo}

Figure~\ref{fig:alg} gives the algorithm as \sys implements it. One pass processes every queue in
turn, as far as its next unsatisfied wait, until no queue can progress. Each queue has a vector clock
$\mathit{VC}[q]$~\cite{mattern1989virtual}. A signal records a copy of its queue's clock, and a wait joins the
recorded clock into its own and keeps a \emph{wait record}. A signal also records a \emph{relay}: for each queue,
the positions of that queue known to be in the signal's first scope along the whole path of semaphores that reached
it, each with the stage filter of a masked signal at its first hop. Barrier structs are appended to a per-queue list. Each
access, and each transition, is handed to \pc{Access}, which compares it with earlier events on the same resource.

\begin{figure*}[tp]
\scriptsize
\begin{minipage}[t]{0.49\textwidth}
\begin{tabbing}
\quad\=\quad\=\quad\=\quad\=\quad\=\kill
\kw{procedure} \pc{Check}$(\tau, \mathit{mode})$ \hfill \emph{Algorithm 1} \\
1\>$\mathit{VC}[q] \gets 0$; $\mathit{Bars}[q], \mathit{Waits}[q] \gets [\,]$; $\mathit{Full}[q] \gets 0$; $\mathit{RA}[q], \mathit{RF}[q] \gets \emptyset$ \quad for all $q$ \\
2\>\kw{repeat} \kw{for} $n$ \kw{in} queues \kw{while} $n$ has a next operation $o$: \\
3\>\>\kw{if} $o = \mathsf{Wait}(x,W)$: \kw{if} $x$ not yet signalled: \kw{break} \\
4\>\>\>$(c, \rho) \gets \mathit{Sig}[x]$; $E \gets \top$ \kw{if} $W{=}\top$ or $n{=}\mathit{host}$ \kw{else} $\ESd(W)$ \\
5\>\>\>\kw{if} $E \setminus \{\mathsf{LT}\} \ne \emptyset$: $\mathit{VC}[n] \gets \max(\mathit{VC}[n], c)$; push $w = (E, |\mathit{Bars}[n]|, \rho)$ onto $\mathit{Waits}[n]$ \\
6\>\>\>\>$\mathit{RA}[n] \gets \mathit{RA}[n] \sqcup \rho$; \kw{if} $E = \top$: $\mathit{RF}[n] \gets \mathit{RF}[n] \sqcup \rho$ and \\
\>\>\>\>\>$\mathit{Full}[n][m] \gets \max(\cdot, t)$ for every unfiltered $(t, \bot) \in \rho[m]$ \\
7\>\>$\mathit{VC}[n][n] \mathrel{+}= 1$; $t \gets \mathit{VC}[n][n]$ \hfill $\triangleright$ $t = p(o)$ \\
8\>\>\kw{if} $o = \mathsf{Sig}(x,S)$: $\mathit{Sig}[x] \gets (\mathit{VC}[n], \pc{Relay}(n, t, S))$ \\
9\>\>\kw{if} $o$ is a barrier struct $B$: append $(t,B)$ to $\mathit{Bars}[n]$; \kw{if} $B$ transitions: \pc{Access}$(t_B)$ \\
10\>\>\kw{if} $o$ is an access: \pc{Access}$(o)$ \\
11\>\kw{until} no queue progressed; report \code{deadlock} for blocked queues \\[3pt]
\kw{function} \pc{Relay}$(n, t, S)$ \quad $\triangleright$ per queue $m$: entries $(t_m, \varphi)$ \\
1\>$\varphi_0 \gets \bot$ \kw{if} $S{=}\top$ \kw{else} $(n, t, |\mathit{Bars}[n]|, \ESs(S))$; $\rho \gets \{n \mapsto (t, \varphi_0)\}$ \\
2\>\kw{if} $S = \top$: \kw{return} $\rho \sqcup \mathit{RA}[n]$ \hfill $\triangleright$ every wait chains into it \\
3\>\kw{if} $\ESs(S) \setminus \{\mathsf{LT}\} \ne \emptyset$: $\rho \gets \rho \sqcup \mathit{RF}[n]$ \\
4\>\kw{for} $w$ \kw{in} $\mathit{Waits}[n]$ with $w.E \ne \top$: \kw{if} $\pc{Chain}(w).X_{\mathit{tot}}$ meets $\ESs(S)$: $\rho \gets \rho \sqcup w.\rho$ \\
5\>\kw{return} $\rho$ \\[3pt]
\kw{function} \pc{Ordered}$(e, u)$ \quad $\triangleright$ $e$ processed before $u$; $m{=}q(e)$, $n{=}q(u)$ \hfill \emph{Algorithm 2} \\
1\>$\mathit{mem} \gets (k(e) = \mathsf{W})$;\ $O \gets R(e) \cap R(u)$ \hfill $\triangleright$ verified configuration$^\dagger$ \\
2\>\kw{if} $m = n = \mathit{host}$: \kw{return} ok \\
3\>\kw{if} $m = n$: $C \gets \pc{Chain}(e)$ advanced over $\mathit{Bars}[n]$ \\
4\>\>\kw{if} $u = t_B$: \kw{if} $C$ reached $B$ last, globally, or locally with $\posof(e){=}\posof(u){\ne}\bot$ \\
\>\>\>\>\kw{or} $e \in \Sc{1}(B)$ (and $e \in \ASc{1}{\ell}(B)\ \forall \ell \in O$ if $\mathit{mem}$; and, if \\
\>\>\>\>$\lambda(B) {\ne} \mathsf{glob}$ and $e$ located, $\lambda(B){=}\mathsf{loc}$ and $\posof(e){=}\posof(u){\ne}\bot$): \kw{return} ok \\
5\>\>\kw{elif} $C.\pc{Reaches}(u, \mathit{mem}, O)$: \kw{return} ok \\
6\>\>\kw{return} ok \kw{if} \pc{ViaWait}$(e,u,\mathit{mem},O)$ \kw{else} \emph{same} \\
7\>\kw{if} $\mathit{VC}[n][m] < p(e)$: \kw{return} \emph{race} \\
8\>\kw{return} ok \kw{if} \pc{ViaWait}$(e,u,\mathit{mem},O)$ \kw{else} \emph{wait} \\[3pt]
\kw{function} \pc{ViaWait}$(e,u,\mathit{mem},O)$ \\
1\>\kw{if} $\mathit{Full}[n][m] \ge p(e)$: \kw{return} true \\
2\>\kw{for} $w$ \kw{in} $\mathit{Waits}[n]$, newest first: \\
3\>\>\kw{if} no $(t, \varphi) \in w.\rho[m]$ has $t \ge p(e)$ and ($\varphi = \bot$ or $\pc{InSignal}(e, \varphi)$): \kw{continue} \\
4\>\>\kw{if} $w.E = \top$: \kw{return} true \hfill $\triangleright$ transitions included \\
5\>\>\kw{if} $u = t_B$: \kw{if} $\pc{Chain}(w).x_{\mathit{last}} = p(u)$: \kw{return} true; \kw{continue} \\
6\>\>\kw{if} $s(u) \in w.E$ \kw{or} $\pc{Chain}(w).\pc{Reaches}(u,\mathit{mem},O)$: \kw{return} true \\
7\>\kw{return} false \\[3pt]
\kw{function} \pc{InSignal}$(e, \varphi = (q, t, k, E_S))$ \quad $\triangleright$ $e$ in the first scope of a masked signal \\
1\>\kw{if} $E_S = \top$ \kw{or} ($e$ not a transition and $s(e) \in E_S$): \kw{return} true \\
2\>$C \gets$ fresh $\pc{Chain}(e)$ advanced over $\mathit{Bars}[q][0..k)$ \\
3\>\kw{return} all reached $M$ and $M^L$ sets meet $E_S$ \kw{if} $k(e){=}\mathsf{W}$ \kw{else} $C.X_{\mathit{tot}} \cup C.X^L_{\mathit{tot}}$ meets $E_S$
\end{tabbing}
\end{minipage}\hfill
\begin{minipage}[t]{0.49\textwidth}
\begin{tabbing}
\quad\=\quad\=\quad\=\quad\=\quad\=\kill
\kw{state} \pc{Chain}$(e)$: global sets $X, M$ and local sets $X^L, M^L$, each \hfill \emph{Algorithm 3} \\
\>as $(\cdot)_o$ for earlier calls and $(\cdot)_c[\cdot]$ per struct of the current call; \\
\>$X_{\mathit{tot}}, X^L_{\mathit{tot}}$; $M_{\mathit{ends}}, M^L_{\mathit{ends}}$; last reached positions; $\mathit{snap}[\cdot]$; $\posof_e$ \\
\>starts at the first barrier at or after $p(e)$ with all sets empty; \\
\>a wait record's chain starts after the wait, seeded $X_o = M_o = X_{\mathit{tot}} = w.E$ \\
\kw{procedure} \pc{Advance} over the next barrier $(t, B)$ of $q(e)$: \\
1\>\kw{if} $B$ sets (resets) event $v$: save (delete) $\mathit{snap}[v]$ of all four sets; \kw{next} \\
2\>\kw{if} $B$ starts a new call: fold $(\cdot)_c$ into $(\cdot)_o$ for all four sets \\
3\>$Y_f \gets Y_o \cup Y_c[B.\mathit{st}]$ for $Y \in \{X, M, X^L, M^L\}$; or $\mathit{snap}[v]$ if $B$ waits on $v$ \\
4\>$d_x \gets e \in \Sc{1}(B)$ (a transition only for its own $B$); $\mathit{lo} \gets d_x \wedge e$ located \\
5\>$d_m \gets d_x \wedge k(e){=}\mathsf{W} \wedge s(e) \in \AS(S_1) \wedge f(e) \in_{\mathsf{W}} A_1 \wedge (r_B{=}\ast \vee (r_B{=}r(e) \wedge R_B \supseteq R(e)))$ \\
\>\>(a transition: $d_m \gets d_x$) \\
6\>\kw{if} $\lambda(B) = \mathsf{glob}$: $g_x \gets d_x \vee (X_f \cup X^L_f)$ meets $\ESs(S_1)$; $l_x \gets$ false \\
\>\kw{else}: $g_x \gets (d_x \wedge \neg\mathit{lo}) \vee X_f$ meets $\ESs(S_1)$; $l_x \gets \mathit{lo} \vee X^L_f$ meets $\ESs(S_1)$ \\
7\>\kw{if} $g_x$: $X_c[B.\mathit{st}], X_{\mathit{tot}} \mathrel{\cup}= \ESd(S_2)$ \hfill $\triangleright$ class $\mathsf{G}$ \\
8\>\kw{if} $l_x$: $X_c[B.\mathit{st}], X_{\mathit{tot}} \mathrel{\cup}= \ESd(S_2) \setminus F$; \hfill $\triangleright$ $\mathsf{N}$ \\
\>\>\kw{if} $\lambda(B) = \mathsf{loc}$: $X^L_c[B.\mathit{st}], X^L_{\mathit{tot}} \mathrel{\cup}= \ESd(S_2) \cap F$ \hfill $\triangleright$ $\mathsf{L}$ \\
9\>the same for $g_m, l_m$ from $d_m$, $M_f$, $M^L_f$, adding $(r_B, R_B, S_2, A_2)$ to $M_{\mathit{ends}}$ \\
\>\>(restricted to stages outside $F$ for $\mathsf{N}$) or to $M^L_{\mathit{ends}}$ (stages in $F$) \\
\kw{function} \pc{Reaches}$(u, \mathit{mem}, O)$: \quad $\mathit{at} \gets \posof_e = \posof(u) \ne \bot$ \\
1\>\kw{if} $\neg\mathit{mem}$: \kw{return} $s(u) \in X_{\mathit{tot}}$ \kw{or} ($\mathit{at}$ \kw{and} $s(u) \in X^L_{\mathit{tot}}$) \\
2\>\kw{return} $\exists (r_B, R_B, \Sigma, A_2)$ in $M_{\mathit{ends}}$, or in $M^L_{\mathit{ends}}$ if $\mathit{at}$: \\
\>\>($r_B{=}\ast$ \kw{or} $r_B{=}r(u) \wedge R_B \supseteq O$) $\wedge$ $s(u) \in \Sigma \cap \AS(S_2) \wedge f(u) \in_{k(u)} A_2$ \\[3pt]
\kw{procedure} \pc{Access}$(u)$ \quad $\triangleright$ per resource: an interval map of segments \hfill \emph{Algorithm 4} \\
\>\>$\triangleright$ each segment $\varsigma$ holds $\mathit{Wr}_\varsigma$, $\mathit{Rd}_\varsigma[\kappa]$, $\mathit{Sh}_\varsigma$ \\
1\>$K \gets \mathit{Wr}_\varsigma$ (+ all $\mathit{Rd}_\varsigma$ if $k(u){=}\mathsf{W}$) for segments $\varsigma \subseteq R(u)$, \\
\>\>+ $\{e : (e,d) \in \mathit{Sh}_\varsigma \text{ live},\ d,u \text{ exempt},\ k(e){=}\mathsf{W} \vee k(u){=}\mathsf{W}\}$ \\
2\>\kw{for} each $e \in K$ once: \kw{if} $e,u$ exempt (E1--E4): $v[e] \gets$ \emph{exempt}; \kw{next} \\
3\>\>$y \gets \pc{Ordered}(e,u)$ \\
4\>\>\kw{if} $y \ne$ ok \kw{and} $(e,u)$ an E5 pair: $v[e] \gets$ \emph{lt}; $K \mathrel{+}= e.\mathit{hid}$ \\
5\>\>\kw{elif} $y \ne$ ok: report $(e,u,y)$; $v[e] \gets y$ \\
6\>\>\kw{elif} $e$ a device write, $u$ a host write, \kw{not} \pc{HostDomain}$(e,u)$: \\
\>\>\>warn $(e,u)$; $v[e] \gets$ \emph{premise} \quad \kw{else} $v[e] \gets$ ok \\
7\>\kw{if} $\mathit{mode} = \mathit{pairwise}$: add $u$ to each $\varsigma$; \kw{return} \\
8\>\kw{for} each segment $\varsigma \subseteq R(u)$: \hfill $\triangleright$ epoch mode \\
9\>\>\kw{if} $k(u) = \mathsf{W}$: \kw{for} $e \in \mathit{Wr}_\varsigma$ with ($v[e]{=}$ok and $u$ not atomic) or \pc{Supersedes}$(u,e)$, \\
\>\>\>\>and $e \in \mathit{Rd}_\varsigma$ with $v[e]{=}$ok and $u$ not atomic: \pc{Drop}$(e,u)$; add $u$ to $\mathit{Wr}_\varsigma$ \\
10\>\>\kw{else} $\kappa \gets (q,s,\text{scope},\text{instance},\text{load/store},\text{atomic},\posof)(u)$; \\
\>\>\>$\mathit{Rd}_\varsigma[\kappa] \gets [u]$, plus the newest read of another command if $u$ has one \\
\pc{Drop}$(e,u)$: remove $e$ from $\varsigma$; \kw{if} $u$ has a command, scope or instance: \\
\>\>add $(e,u)$ to $\mathit{Sh}_\varsigma$; \kw{if} $u$ is a transition: $u.\mathit{hid} \mathrel{+}= e$ \hfill $\triangleright$ R-LT \\
\pc{Supersedes}$(u,e)$: same queue, stage, flag, scope, instance, load/store \\
\>marker, atomic mode and token; both writes, neither a transition \\
\>nor in a barrier call, and not one command \\
\pc{HostDomain}$(e,u)$: $C \gets$ fresh \pc{Chain}$(e)$ over $\mathit{Bars}[q(e)]$ up to $\mathit{Full}[\mathit{host}][q(e)]$; \\
\>\kw{return} $C.\pc{Reaches}(h, \mathsf{true}, O)$ for a host access $h$ of $r(e)$ with \\
\>flag \code{HOST\_READ} or \code{HOST\_WRITE} \hfill $\triangleright$ Def.~\ref{def:hostdom}
\end{tabbing}
\end{minipage}
\caption{The checking algorithm as implemented in \sys (\code{covenant/timing.py}), verified configuration. ``Meets''
is intersection of stage sets, where $\top$ meets every set that has a stage other than $\mathsf{LT}$. A relay $\rho$
maps each queue $m$ to entries $(t, \varphi)$: the operations of $m$ at positions $\le t$ are in the first scope of
the signal along the whole path, all of them if $\varphi = \bot$, or only those that pass \pc{InSignal} for the
masked signal $\varphi$ at the first hop. $\sqcup$ merges relays, keeping per queue the largest unfiltered entry and
at most four newer filtered ones; dropping an entry only removes orderings. $\mathit{RA}[n]$ and $\mathit{RF}[n]$
merge the relays of all waits and of the full waits on $n$. A shadow entry is live while its dominator's command,
scope or instance is still current on its queue. In Algorithm~3, the global sets hold links of classes $\mathsf{G}$
and $\mathsf{N}$ and the local sets links of class $\mathsf{L}$ (Definition~\ref{def:spec}); a local link orders only
an access with the source's token $\posof_e$. \pc{InSignal} tests the union of global and local sets, since a signal
is an ordinary link. $^\dagger$Under the opt-in heuristic R1, $\mathit{mem}$ is false when $u$ is a transition from
\code{UNDEFINED}; that mode is outside the theorems.}
\label{fig:alg}
\end{figure*}

\paragraph{Invariants}
Four invariants carry the proof, and they are the reason for each data structure of Figure~\ref{fig:alg}.
\begin{enumerate}
\item \emph{Processing order is a linearization.} Operations are processed queue by queue, each wait only after its
  signal, so every chain of Definition~\ref{def:spec} runs forward in processing order (Lemma~\ref{lem:lin}). An
  earlier event is therefore always available when a later one is checked, and a queue whose clock entry for $m$ is
  below $p(e)$ has no chain from $e$ at all (Algorithm~2, line~7).
\item \emph{The chain state is exactly the reachable links.} For a source event $e$, \pc{Chain}$(e)$ walks the barriers
  after $e$ on its queue once, incrementally, and keeps the union of the second stage sets of the links it has
  reached, split by class: $X$ and $M$ hold the $\mathsf{G}$ and $\mathsf{N}$ links (orderings valid for every access
  of those stages) and $X^L$, $M^L$ the $\mathsf{L}$ links (valid only at $e$'s own framebuffer position). The structs
  of the current call are kept apart, because structs of one call never chain into one another, and an event wait
  continues only from what had been reached at its event's set. The memory sets admit a first link only if it holds
  all of $R(e)$ in its first access scope (Lemma~\ref{lem:chain}).
\item \emph{A relay entry is a chain into the signal.} An entry $(t,\varphi)$ for queue $m$ in a signal's relay
  means that every operation of $m$ up to $t$ that passes the stage filter $\varphi$ reaches the signal by a chain
  (Lemma~\ref{lem:clock}). A wait record keeps that relay, so a cross-queue question becomes a same-queue question on
  the waiting queue: does the wait's own chain reach $u$ (Algorithm~2, \pc{ViaWait})?
\item \emph{Every dropped entry has a live stand-in.} In epoch mode an entry is removed only when a later access
  orders everything it would order: a write ordered after it, a write with the same stage, flag, scope, instance,
  atomic mode and token, or a read with the same key (including the token). Where the stand-in may be exempt from a
  later access that the dropped entry is not (same command, scope or instance), the entry is kept in a shadow list
  (Lemma~\ref{lem:fast}).
\end{enumerate}
In \emph{pairwise} mode no entry is dropped, so every earlier conflicting event is compared with every later one. In
\emph{epoch} mode, in the manner of FastTrack~\cite{flanagan2009fasttrack}, invariant~4 keeps the candidate lists
short. A transition keeps the entries it dropped, because under R-LT it does not stand in for them against a later
transition. Theorem~\ref{thm:sound} is stated for the pairwise mode, which the optimizer also uses to re-check its
final plans (\S\ref{sec:elide}). The epoch mode is related to it as follows.

\begin{lemma}[Modes]\label{lem:modes}
\pc{Ordered}$(e,u)$ does not depend on the mode, so $\Hep(\tau) \subseteq \Hpw(\tau)$.
\end{lemma}

\begin{lemma}[Fast mode]\label{lem:fast}
Under the premises of Theorem~\ref{thm:sound} and W2, every race $\{a,b\} \in \Hspec(\tau)$ has $a$ or $b$ in
some pair of $\Hep(\tau)$ (criterion E). In particular, $\Hep(\tau) = \emptyset$ implies $\Hspec(\tau) = \emptyset$.
\end{lemma}

The proof (Appendix~\ref{app:fast}) shows that every drop keeps a stand-in: a later access $x$ that races the dropped
entry $e$ either is compared with an access $v$ that stands in for $e$, where an \emph{ok} for $(v,x)$ would order $e$
before $x$ by transitivity of Definition~\ref{def:spec} through $v$, or finds $e$ in the shadow list. Two remarks
qualify the lemma. First, it relates $\Hep$ to $\Hspec$, not to $\Hpw$. An empty epoch result does not imply an empty
pairwise result: when an entry is dropped piecewise by stand-ins with smaller regions, the epoch mode can be more
precise than the pairwise mode, which tests regions as a whole. Second, two stand-in rules are necessary. An access
with an atomic mode never stands in for what it is ordered after: after a fill of a counter, a barrier from transfer
to compute, an atomic add in a compute shader and an atomic add in a fragment shader on one queue, the fragment add is
exempt from the compute add (E4) but races the fill. And a read stands in only for a read with the same token: after a
sampled read and an input-attachment read of one image in a fragment shader, a \code{BY\_REGION} dependency and a
color write at the same position, the write is ordered after the input-attachment read but races the sampled read.
Both frames are regression tests (Appendix~\ref{app:counter}).

\paragraph{Diagnostics}
\begin{sloppypar}
Diagnostics include \code{timing.missing-barrier}, \code{barrier-scope}, \code{cross-queue-race},
\code{wait-stage-scope}, \code{signal-scope}, \code{wrong-layout}, \code{stale-barrier}, \code{ownership-missing},
\code{free-while-in-use} and \code{deadlock}, and the warnings \code{unmodelled-command} and
\code{premise-host-domain}. Each same-queue report names the
scope test that failed, in the specification's terms: the first synchronization scope (execution), the first access
scope (availability), the second access scope (visibility), or a broken chain. A cross-queue report through a relay
whose stage mask cuts the chain is \code{signal-stage-scope}.
\end{sloppypar}

\paragraph{Protocol checks are not race checks}
\label{sec:protocol}
Some valid-usage rules concern synchronization objects, not memory, and Theorem~\ref{thm:sound} says nothing about
them. \sys checks them along the recorded host order, and each is justified by the rule it cites, not by
Definition~\ref{def:spec}.
\begin{sloppypar}
\begin{itemize}
\item \code{swapchain-overacquire}. Let $S$ be the swapchain's image count and $M$ the surface's
  \code{minImageCount}. The specification says that \code{vkAcquireNextImageKHR} should not be called while the
  application holds more than $S-M$ images. It guarantees that an acquire returns in finite time only when at most
  $S-M$ are held, and it forbids an infinite timeout when that forward progress cannot be guaranteed
  (\code{VUID-vkAcquireNextImageKHR-}\allowbreak\code{surface-07783}). The check counts acquires minus presents in host order and
  reports an acquire made while more than $S-M$ images are held. With an infinite timeout this is a
  valid-usage violation that can block forever; with a finite one, a call the specification advises against. The
  count is exact for the recorded order and says nothing about other driver responses.
\item \code{present-semaphore-reuse}. A binary semaphore must be unsignaled when a signal operation on it executes
  (\code{VUID-vkQueueSubmit-}\allowbreak\code{pSignalSemaphores-00067}, \code{VUID-vkQueueSubmit2-}\allowbreak\code{semaphore-03868},
  \code{VUID-vkAcquireNextImageKHR-}\allowbreak\code{semaphore-01286}). A present unsignals its wait semaphores at a time that no
  submission fence bounds. The specification offers two ways to know that it has happened: a present fence
  (\code{VK\_KHR\_swapchain\_maintenance1}), or acquiring the same image again and waiting on that acquire's semaphore
  or fence. The justified rule is therefore an ordering obligation: a re-signal must happen after the signal of a
  later acquisition of the same image. The check reports a re-signal made before the image is acquired again. It
  approximates ``happens after the acquisition's signal'' by ``follows the acquire call in host order''. It therefore
  accepts a re-signal from a batch that does not wait on that acquisition, and it does not model present fences.
\item The in-use rules (a command buffer, fence or acquire semaphore reused while pending:
  \code{VUID-vkBeginCommandBuffer-}\allowbreak\code{commandBuffer-00049}, \code{VUID-vkResetFences-}\allowbreak\code{pFences-01123},
  \code{VUID-vkAcquireNextImageKHR-}\allowbreak\code{semaphore-01779}) and \code{free-while-in-use}
  (\code{VUID-vkFreeDescriptorSets-}\allowbreak\code{pDescriptorSets-00309}) are encoded as accesses to pseudo-resources and checked by
  the same order. Theorem~\ref{thm:sound} applies to them formally, but their meaning comes from these rules, not
  from the memory model. A binary semaphore signalled while a signal is pending, a timeline value that does not
  increase, and a wait that is never satisfied are reported by the trace translation.
\end{itemize}
\end{sloppypar}

\paragraph{Report grouping by checked repair}
Raw diagnostics are produced per accessing command and per frame, which is far too many to read. \sys reduces them
in three steps:
\begin{enumerate}
\item It canonicalizes each diagnostic to a site key: code, resource role, command site and barrier site.
\item For each group it synthesizes the narrowest barrier or mask repair and re-runs the checker in the pairwise
  mode. A repair is \emph{checked} when the group's hazards disappear and no racing pair appears that was not
  reported before, with pairs compared by event identity (call, command and access index), the same rule as for
  barrier transformations (\S\ref{sec:elide}). Groups that one repair eliminates are merged.
\item It labels as technicalities the hazards on discarded content and the ownership issues that only matter with
  a single queue family.
\end{enumerate}
The result groups diagnostics by site and by checked repair. It is not one report per mistake: one root cause can
still produce several groups, for example once per recreated swapchain, and we measure how often in
\S\ref{sec:eval-fresh}. Grouping runs the checker as an oracle within a budget of hazard pairs and memory. When a run
exceeds it, grouping falls back to site groups without repairs and marks the result \emph{incomplete}; every
diagnostic is still a member of exactly one group, and an incomplete result is never presented as clean. Lemma~\ref{lem:repair} below shows that the repairs \sys synthesizes cannot add a race under
the specification. It says nothing about the checker's reports, which over-approximate, so the event-pair comparison
remains a safeguard that may reject a repair whose new reported pair is not a race.

\begin{lemma}[Repair monotonicity]\label{lem:repair}
Let $\tau_r$ be $\tau$ with a repair $r$ applied, where $r$ either widens the stage and access masks of one
dependency (an explicit barrier struct, a subpass dependency or the implicit external dependency of a render pass,
whose masks the explicit one then contains) without changing its layouts or regions, or inserts a barrier without
a layout transition. Then the events of $\tau_r$ are those of $\tau$, identified as above, and
$\Hspec(\tau_r) \subseteq \Hspec(\tau)$.
\end{lemma}

The proof is in Appendix~\ref{app:repair}. A source fix must keep the lemma's premise: replacing an implicit external
dependency by an explicit one with \emph{smaller} masks is not a widening. As a test, every candidate repair the
grouper synthesizes on 300 random frames (4,791 repairs) was applied and checked against the executable reference
semantics; none added a racing pair.
